\documentclass[10pt,a4paper]{article}

\usepackage{amsmath}
\usepackage{amssymb}
\usepackage{amsfonts}
\usepackage{enumerate}
\usepackage[top=2.5cm, bottom=2.5cm, left=2.5cm, right=2.5cm]{geometry}
\usepackage{multicol}
\usepackage{parskip}
\usepackage{titlesec}
\usepackage{fancyhdr}
\usepackage{xcolor}

\pagestyle{fancy}
\fancyhf{}
\rhead{Chapter Zero: Mathematical Prerequisites}
\lhead{MCQ Practice Set}
\cfoot{\thepage}

\titleformat{\section}{\large\bfseries}{}{0em}{}[\titlerule]

\newcommand{\topic}[1]{\section*{\large\textsc{#1}}}

\begin{document}

\begin{center}
   % {\LARGE \textbf{Chapter Zero: Mathematical Prerequisites}}\\[0.4em]
    {\large MCQ Question Bank}\\[0.2em]
    %{\normalsize Topics: Logarithm $\cdot$ Trigonometry $\cdot$ Vectors $\cdot$ Graphs $\cdot$ Differentiation $\cdot$ Integration}\\[0.2em]
   % {\small Difficulty: Easy $\longrightarrow$ NEST Level}
\end{center}

\vspace{0.5em}
\hrule
\vspace{1em}

%% ============================================================
\topic{Section 1: Indices and Logarithms}
%% ============================================================

\begin{enumerate}[1.]

\item The value of $\log_2 128$ is \hfill \textbf{(C)}
\begin{enumerate}[(A)]
    \item 5
    \item 6
    \item 7
    \item 8
\end{enumerate}

\item $\log_a 1$ is equal to \hfill \textbf{(B)}
\begin{enumerate}[(A)]
    \item 1
    \item 0
    \item $a$
    \item undefined
\end{enumerate}

\item If $\log_{10} 2 = 0.3010$, then $\log_{10} 50$ equals \hfill \textbf{(A)}
\begin{enumerate}[(A)]
    \item $1.6990$
    \item $1.7782$
    \item $1.5000$
    \item $0.6990$
\end{enumerate}

\item Which of the following is equal to $\log_a\!\left(\dfrac{m}{n}\right)$? \hfill \textbf{(C)}
\begin{enumerate}[(A)]
    \item $\log_a m \cdot \log_a n$
    \item $\log_a m + \log_a n$
    \item $\log_a m - \log_a n$
    \item $\dfrac{\log_a m}{\log_a n}$
\end{enumerate}

\item The number of digits in $2^{10}$ (given $\log_{10} 2 = 0.3010$) is \hfill \textbf{(B)}
\begin{enumerate}[(A)]
    \item 3
    \item 4
    \item 5
    \item 6
\end{enumerate}

\item $\log_a b \times \log_b a$ equals \hfill \textbf{(D)}
\begin{enumerate}[(A)]
    \item $\log_a b$
    \item $\log_b a$
    \item 0
    \item 1
\end{enumerate}

\item If $\log(x+y) = \log x + \log y$, then $x$ expressed in terms of $y$ is \hfill \textbf{(B)}
\begin{enumerate}[(A)]
    \item $\dfrac{y}{1+y}$
    \item $\dfrac{y}{y-1}$
    \item $\dfrac{1}{y-1}$
    \item $\dfrac{y^2}{y-1}$
\end{enumerate}

\item $\log_4 8 + \log_8 32 - \log_{16} 8$ equals \hfill \textbf{(C)}
\begin{enumerate}[(A)]
    \item $2$
    \item $\dfrac{5}{3}$
    \item $\dfrac{29}{12}$
    \item $\dfrac{7}{4}$
\end{enumerate}

\item Given $\log 2 = a$ and $\log 3 = b$, the value of $\log_{12} 6\sqrt{2}$ is \hfill \textbf{(A)}
\begin{enumerate}[(A)]
    \item $\dfrac{3a+2b}{4a+2b}$
    \item $\dfrac{2a+b}{3a+2b}$
    \item $\dfrac{a+b}{2a+b}$
    \item $\dfrac{3a+b}{2a+2b}$
\end{enumerate}

\item The value of $\log_a(a^x)$ is \hfill \textbf{(B)}
\begin{enumerate}[(A)]
    \item $a^x$
    \item $x$
    \item $x \cdot a$
    \item $\log x$
\end{enumerate}

\item If $2^x = 3^y = 6^{-z}$, then $\dfrac{1}{x} + \dfrac{1}{y} + \dfrac{1}{z}$ equals \hfill \textbf{(C)}
\begin{enumerate}[(A)]
    \item 1
    \item $-1$
    \item 0
    \item $\dfrac{1}{6}$
\end{enumerate}

\item The expression $\log a^2 + \log b^3 - \log c^4$ simplifies to \hfill \textbf{(D)}
\begin{enumerate}[(A)]
    \item $\log\!\left(\dfrac{a^2 b^3}{c^4}\right)^2$
    \item $6\log(abc)$
    \item $\log(a^2 - b^3 + c^4)$
    \item $\log\!\left(\dfrac{a^2 b^3}{c^4}\right)$
\end{enumerate}

\item If $a^2 + b^2 = 23ab$, then $\log\!\left(\dfrac{a+b}{5}\right)$ equals \hfill \textbf{(A)}
\begin{enumerate}[(A)]
    \item $\dfrac{1}{2}(\log a + \log b)$
    \item $\log a + \log b$
    \item $\dfrac{1}{2}\log(a-b)$
    \item $2(\log a + \log b)$
\end{enumerate}

\item The value of $\log_a b \times \log_b c \times \log_c a$ is \hfill \textbf{(B)}
\begin{enumerate}[(A)]
    \item 0
    \item 1
    \item $\log(abc)$
    \item $abc$
\end{enumerate}

\item $\ln(x)$ in terms of $\log_{10}(x)$ is approximately \hfill \textbf{(C)}
\begin{enumerate}[(A)]
    \item $0.4343 \cdot \log_{10}(x)$
    \item $10 \cdot \log_{10}(x)$
    \item $2.303 \cdot \log_{10}(x)$
    \item $e \cdot \log_{10}(x)$
\end{enumerate}

\item If $4^{(x+1)} + 4^{(1-x)} = 10$, the values of $x$ are \hfill \textbf{(D)}
\begin{enumerate}[(A)]
    \item $x = 0$ or $x = 1$
    \item $x = \dfrac{1}{4}$ or $x = \dfrac{3}{4}$
    \item $x = -1$ or $x = 1$
    \item $x = -\dfrac{1}{2}$ or $x = \dfrac{1}{2}$
\end{enumerate}

\item $\log_a\!\left(\sqrt{x^2+1}+x\right) + \log_a\!\left(\sqrt{x^2+1}-x\right)$ is equal to \hfill \textbf{(B)}
\begin{enumerate}[(A)]
    \item $2\log_a x$
    \item $0$
    \item $1$
    \item $\log_a 2$
\end{enumerate}

\item Which of the following is NOT a valid domain condition for $\log_a b$? \hfill \textbf{(C)}
\begin{enumerate}[(A)]
    \item $a > 0$
    \item $b > 0$
    \item $a = 1$ is allowed
    \item $a \neq 1$
\end{enumerate}

\item If $\log_2 3 = a$, then $\log_8 48$ equals \hfill \textbf{(A)}
\begin{enumerate}[(A)]
    \item $\dfrac{4+a}{3}$
    \item $\dfrac{3+a}{4}$
    \item $\dfrac{a+2}{3}$
    \item $4+a$
\end{enumerate}

\item The sum $\dfrac{\log a}{\log ab} + \dfrac{\log b}{\log ab}$ evaluates to \hfill \textbf{(D)}
\begin{enumerate}[(A)]
    \item $\log(ab)$
    \item $\log a + \log b$
    \item 0
    \item 1
\end{enumerate}

\item If $x = 2^{1/3} + 2^{-1/3}$, then $2x^3 - 6x$ equals \hfill \textbf{(C)}
\begin{enumerate}[(A)]
    \item 2
    \item 3
    \item 5
    \item 7
\end{enumerate}

\item The simplified form of $\dfrac{2^5 \times 4^3}{8^2 \times 16}$ is \hfill \textbf{(B)}
\begin{enumerate}[(A)]
    \item 1
    \item 2
    \item 4
    \item 8
\end{enumerate}

\end{enumerate}

%% ============================================================
\newpage
\topic{Section 2: Trigonometry}
%% ============================================================

\begin{enumerate}[1.]

\item The value of $\sin 30^\circ \cdot \cos 60^\circ + \cos 30^\circ \cdot \sin 60^\circ$ is \hfill \textbf{(D)}
\begin{enumerate}[(A)]
    \item $0$
    \item $\dfrac{\sqrt{3}}{2}$
    \item $\dfrac{1}{2}$
    \item $1$
\end{enumerate}

\item $\cos(-\theta)$ equals \hfill \textbf{(A)}
\begin{enumerate}[(A)]
    \item $\cos\theta$
    \item $-\cos\theta$
    \item $\sin\theta$
    \item $-\sin\theta$
\end{enumerate}

\item The exact value of $\cos 75^\circ$ is \hfill \textbf{(C)}
\begin{enumerate}[(A)]
    \item $\dfrac{\sqrt{3}+1}{2\sqrt{2}}$
    \item $\dfrac{\sqrt{6}+\sqrt{2}}{4}$
    \item $\dfrac{\sqrt{6}-\sqrt{2}}{4}$
    \item $\dfrac{\sqrt{3}-1}{2}$
\end{enumerate}

\item From the Pythagorean identity $\sin^2\theta + \cos^2\theta = 1$, dividing by $\cos^2\theta$ gives \hfill \textbf{(B)}
\begin{enumerate}[(A)]
    \item $1 + \cot^2\theta = \csc^2\theta$
    \item $1 + \tan^2\theta = \sec^2\theta$
    \item $\tan^2\theta - \sec^2\theta = 1$
    \item $\sin^2\theta = 1 + \cos^2\theta$
\end{enumerate}

\item If $\tan\theta = \dfrac{3}{4}$ and $\theta$ is in the third quadrant, then $\cos\theta$ equals \hfill \textbf{(C)}
\begin{enumerate}[(A)]
    \item $\dfrac{4}{5}$
    \item $\dfrac{3}{5}$
    \item $-\dfrac{4}{5}$
    \item $-\dfrac{3}{5}$
\end{enumerate}

\item For small angle $\theta$ in radians, $\cos\theta \approx$ \hfill \textbf{(A)}
\begin{enumerate}[(A)]
    \item $1$
    \item $\theta$
    \item $1 - \theta$
    \item $\theta^2$
\end{enumerate}

\item The value of $\sin 20^\circ \cdot \sin 40^\circ \cdot \sin 80^\circ$ is \hfill \textbf{(D)}
\begin{enumerate}[(A)]
    \item $\dfrac{1}{4}$
    \item $\dfrac{\sqrt{3}}{4}$
    \item $\dfrac{1}{8}$
    \item $\dfrac{\sqrt{3}}{8}$
\end{enumerate}

\item $\sin A + \sin B$ equals \hfill \textbf{(B)}
\begin{enumerate}[(A)]
    \item $2\cos\!\left(\dfrac{A+B}{2}\right)\sin\!\left(\dfrac{A-B}{2}\right)$
    \item $2\sin\!\left(\dfrac{A+B}{2}\right)\cos\!\left(\dfrac{A-B}{2}\right)$
    \item $2\sin\!\left(\dfrac{A+B}{2}\right)\sin\!\left(\dfrac{A-B}{2}\right)$
    \item $2\cos\!\left(\dfrac{A+B}{2}\right)\cos\!\left(\dfrac{A-B}{2}\right)$
\end{enumerate}

\item The value of $\sin(A+B)\sin(A-B)$ simplifies to \hfill \textbf{(A)}
\begin{enumerate}[(A)]
    \item $\sin^2 A - \sin^2 B$
    \item $\cos^2 A - \cos^2 B$
    \item $\sin^2 A + \sin^2 B$
    \item $\cos^2 B - \cos^2 A$
\end{enumerate}

\item The double angle formula $\cos 2\theta$ can be written as \hfill \textbf{(D)}
\begin{enumerate}[(A)]
    \item $2\sin\theta\cos\theta$
    \item $1 - 2\cos^2\theta$
    \item $\dfrac{1-\tan^2\theta}{1+\tan^2\theta}$ only
    \item $\cos^2\theta - \sin^2\theta = 1 - 2\sin^2\theta = 2\cos^2\theta - 1$
\end{enumerate}

\item If $\tan(\theta/2) = t$, then $\sin\theta$ equals \hfill \textbf{(B)}
\begin{enumerate}[(A)]
    \item $\dfrac{1-t^2}{1+t^2}$
    \item $\dfrac{2t}{1+t^2}$
    \item $\dfrac{2t}{1-t^2}$
    \item $\dfrac{t}{1+t^2}$
\end{enumerate}

\item In a triangle $ABC$, $\tan A + \tan B + \tan C$ equals \hfill \textbf{(C)}
\begin{enumerate}[(A)]
    \item $0$
    \item $\tan A + \tan B - \tan C$
    \item $\tan A \cdot \tan B \cdot \tan C$
    \item $1$
\end{enumerate}

\item All $\theta \in [0^\circ, 360^\circ)$ satisfying $2\cos^2\theta - 3\cos\theta + 1 = 0$ are \hfill \textbf{(A)}
\begin{enumerate}[(A)]
    \item $0^\circ,\, 60^\circ,\, 300^\circ$
    \item $0^\circ,\, 90^\circ,\, 270^\circ$
    \item $60^\circ,\, 120^\circ,\, 300^\circ$
    \item $0^\circ,\, 120^\circ,\, 240^\circ$
\end{enumerate}

\item The value of $\tan\!\left(\dfrac{13\pi}{12}\right)$ is \hfill \textbf{(D)}
\begin{enumerate}[(A)]
    \item $\sqrt{3}$
    \item $2+\sqrt{3}$
    \item $1-\sqrt{3}$
    \item $2-\sqrt{3}$
\end{enumerate}

\item $(\sec\theta - \tan\theta)^2$ simplifies to \hfill \textbf{(B)}
\begin{enumerate}[(A)]
    \item $\dfrac{1+\sin\theta}{1-\sin\theta}$
    \item $\dfrac{1-\sin\theta}{1+\sin\theta}$
    \item $\dfrac{1-\cos\theta}{1+\cos\theta}$
    \item $\dfrac{1+\cos\theta}{1-\cos\theta}$
\end{enumerate}

\item In triangle $ABC$, the sine rule states that \hfill \textbf{(C)}
\begin{enumerate}[(A)]
    \item $a^2 = b^2 + c^2 - 2bc\cos A$
    \item $\dfrac{\sin A}{a} = \dfrac{\sin B}{b} = R$
    \item $\dfrac{\sin A}{a} = \dfrac{\sin B}{b} = \dfrac{\sin C}{c}$
    \item $\dfrac{a}{\sin A} = 2R^2$
\end{enumerate}

\item $\sin^2\!\left(\dfrac{\theta}{2}\right)$ in terms of $\cos\theta$ is \hfill \textbf{(A)}
\begin{enumerate}[(A)]
    \item $\dfrac{1-\cos\theta}{2}$
    \item $\dfrac{1+\cos\theta}{2}$
    \item $\dfrac{\cos\theta - 1}{2}$
    \item $1 - \cos^2\theta$
\end{enumerate}

\item The expression $(\sin\theta + \csc\theta)^2 + (\cos\theta + \sec\theta)^2$ equals \hfill \textbf{(D)}
\begin{enumerate}[(A)]
    \item $\tan^2\theta + \cot^2\theta + 5$
    \item $\tan^2\theta + \cot^2\theta + 9$
    \item $\sec^2\theta + \csc^2\theta$
    \item $\tan^2\theta + \cot^2\theta + 7$
\end{enumerate}

\item The solution set of $\sin x + \sin 3x + \sin 5x = 0$ in $[0^\circ, 360^\circ)$ has how many elements? \hfill \textbf{(C)}
\begin{enumerate}[(A)]
    \item 4
    \item 5
    \item 6
    \item 7
\end{enumerate}

\item If $A + B + C = \pi$, then $\sin 2A + \sin 2B + \sin 2C$ equals \hfill \textbf{(B)}
\begin{enumerate}[(A)]
    \item $2\sin A\sin B\sin C$
    \item $4\sin A\sin B\sin C$
    \item $4\cos A\cos B\cos C$
    \item $0$
\end{enumerate}

\item In the simple pendulum approximation for small oscillations, the angular frequency $\omega$ for a pendulum of length $L$ is \hfill \textbf{(A)}
\begin{enumerate}[(A)]
    \item $\sqrt{g/L}$
    \item $\sqrt{L/g}$
    \item $g/L$
    \item $2\pi\sqrt{L/g}$
\end{enumerate}

\item The value of $\dfrac{\sin 7A + \sin A}{\cos 7A + \cos A}$ is \hfill \textbf{(C)}
\begin{enumerate}[(A)]
    \item $\tan 3A$
    \item $\tan 7A$
    \item $\tan 4A$
    \item $\cot 4A$
\end{enumerate}

\end{enumerate}

%% ============================================================
\newpage
\topic{Section 3: Vectors}
%% ============================================================

\begin{enumerate}[1.]

\item A unit vector has magnitude equal to \hfill \textbf{(B)}
\begin{enumerate}[(A)]
    \item 0
    \item 1
    \item depends on the vector
    \item $\sqrt{3}$
\end{enumerate}

\item The magnitude of $\vec{A} = 3\hat{i} - 4\hat{j}$ is \hfill \textbf{(C)}
\begin{enumerate}[(A)]
    \item 1
    \item 3
    \item 5
    \item 7
\end{enumerate}

\item The dot product $\vec{A} \cdot \vec{B} = 0$ implies \hfill \textbf{(D)}
\begin{enumerate}[(A)]
    \item $\vec{A} = 0$ or $\vec{B} = 0$
    \item $\vec{A}$ is parallel to $\vec{B}$
    \item $|\vec{A}| = |\vec{B}|$
    \item $\vec{A}$ is perpendicular to $\vec{B}$ (or one is zero)
\end{enumerate}

\item The cross product $\hat{i} \times \hat{j}$ equals \hfill \textbf{(A)}
\begin{enumerate}[(A)]
    \item $\hat{k}$
    \item $-\hat{k}$
    \item $\hat{i}$
    \item $0$
\end{enumerate}

\item What is the dot product of two vectors of magnitudes 3 and 5 if the angle between them is $60^\circ$? \hfill \textbf{(B)}
\begin{enumerate}[(A)]
    \item $5.2$
    \item $7.5$
    \item $8.4$
    \item $15$
\end{enumerate}

\item A vector is NOT changed when \hfill \textbf{(D)}
\begin{enumerate}[(A)]
    \item it is rotated through an arbitrary angle
    \item it is multiplied by a scalar other than 1
    \item it is cross-multiplied by another vector
    \item it is displaced parallel to itself
\end{enumerate}

\item The unit vector in the direction of $\vec{A} = 3\hat{i} - 4\hat{j}$ is \hfill \textbf{(C)}
\begin{enumerate}[(A)]
    \item $3\hat{i} - 4\hat{j}$
    \item $\dfrac{3\hat{i} - 4\hat{j}}{3}$
    \item $\dfrac{3\hat{i} - 4\hat{j}}{5}$
    \item $\dfrac{3\hat{i} - 4\hat{j}}{7}$
\end{enumerate}

\item If $|\vec{A} + \vec{B}| = |\vec{A} - \vec{B}|$, then the angle between $\vec{A}$ and $\vec{B}$ is \hfill \textbf{(B)}
\begin{enumerate}[(A)]
    \item $0^\circ$
    \item $90^\circ$
    \item $60^\circ$
    \item $180^\circ$
\end{enumerate}

\item Three vectors $\vec{A}$, $\vec{B}$, $\vec{C}$ satisfy $\vec{A}+\vec{B}+\vec{C}=\vec{0}$ with $|\vec{A}|=3$, $|\vec{B}|=5$, $|\vec{C}|=7$. The angle between $\vec{A}$ and $\vec{B}$ is \hfill \textbf{(D)}
\begin{enumerate}[(A)]
    \item $30^\circ$
    \item $45^\circ$
    \item $90^\circ$
    \item $60^\circ$
\end{enumerate}

\item A vector makes angles $60^\circ$, $45^\circ$, and $\alpha$ with the $x$, $y$, $z$ axes. A possible value of $\alpha$ is \hfill \textbf{(C)}
\begin{enumerate}[(A)]
    \item $30^\circ$
    \item $45^\circ$
    \item $60^\circ$
    \item $90^\circ$
\end{enumerate}

\item The cross product $\vec{A} \times \vec{B}$ is anti-commutative, meaning \hfill \textbf{(A)}
\begin{enumerate}[(A)]
    \item $\vec{A} \times \vec{B} = -\vec{B} \times \vec{A}$
    \item $\vec{A} \times \vec{B} = \vec{B} \times \vec{A}$
    \item $\vec{A} \times \vec{B} = 0$
    \item $|\vec{A} \times \vec{B}| = \vec{A} \cdot \vec{B}$
\end{enumerate}

\item The area of the parallelogram formed by $\vec{A} = \hat{i}+2\hat{j}+3\hat{k}$ and $\vec{B} = -3\hat{i}-2\hat{j}+\hat{k}$ is \hfill \textbf{(D)}
\begin{enumerate}[(A)]
    \item $3\sqrt{5}$
    \item $\sqrt{180}$
    \item $4\sqrt{5}$
    \item $6\sqrt{5}$
\end{enumerate}

\item If a unit vector is $0.5\hat{i} + 0.8\hat{j} + c\hat{k}$, then $c$ equals \hfill \textbf{(B)}
\begin{enumerate}[(A)]
    \item $\sqrt{0.39}$
    \item $\sqrt{0.11}$
    \item $\sqrt{0.01}$
    \item $0.39$
\end{enumerate}

\item Given $\vec{A}\cdot\vec{B} = \vec{A}\cdot\vec{C}$ and $\vec{A} \neq \vec{0}$, does $\vec{B} = \vec{C}$ necessarily? \hfill \textbf{(C)}
\begin{enumerate}[(A)]
    \item Yes, always
    \item Yes, only if $|\vec{A}|=1$
    \item No; $\vec{B}-\vec{C}$ could be perpendicular to $\vec{A}$
    \item Yes, if $\vec{A}$, $\vec{B}$, $\vec{C}$ are coplanar
\end{enumerate}

\item The torque $\vec{\tau}$ due to a force $\vec{F} = (3\hat{i}+4\hat{j}-\hat{k})$ N at position $\vec{r} = (2\hat{i}-\hat{j}+\hat{k})$ m is \hfill \textbf{(A)}
\begin{enumerate}[(A)]
    \item $-3\hat{i}+5\hat{j}+11\hat{k}$ N$\cdot$m
    \item $3\hat{i}-5\hat{j}+11\hat{k}$ N$\cdot$m
    \item $-3\hat{i}-5\hat{j}+11\hat{k}$ N$\cdot$m
    \item $3\hat{i}+5\hat{j}-11\hat{k}$ N$\cdot$m
\end{enumerate}

\item Two forces of 5 N and 12 N act at right angles to each other. The magnitude of their resultant is \hfill \textbf{(C)}
\begin{enumerate}[(A)]
    \item 7 N
    \item 17 N
    \item 13 N
    \item 60 N
\end{enumerate}

\item The identity $|\vec{A}+\vec{B}|^2 + |\vec{A}-\vec{B}|^2$ equals \hfill \textbf{(B)}
\begin{enumerate}[(A)]
    \item $|\vec{A}|^2 + |\vec{B}|^2$
    \item $2(|\vec{A}|^2 + |\vec{B}|^2)$
    \item $4\,\vec{A}\cdot\vec{B}$
    \item $2\,\vec{A}\cdot\vec{B}$
\end{enumerate}

\item If $\vec{A}\times\vec{B} = \vec{C}\times\vec{D}$ and $\vec{A}\times\vec{C} = \vec{B}\times\vec{D}$, then $(\vec{A}-\vec{D})$ is \hfill \textbf{(D)}
\begin{enumerate}[(A)]
    \item perpendicular to $(\vec{B}-\vec{C})$
    \item anti-parallel to $(\vec{B}-\vec{C})$
    \item equal to $(\vec{B}-\vec{C})$
    \item parallel to $(\vec{B}-\vec{C})$
\end{enumerate}

\item The direction cosines $\cos l$, $\cos m$, $\cos n$ of any vector satisfy \hfill \textbf{(A)}
\begin{enumerate}[(A)]
    \item $\cos^2 l + \cos^2 m + \cos^2 n = 1$
    \item $\cos l + \cos m + \cos n = 1$
    \item $\cos^2 l + \cos^2 m + \cos^2 n = 3$
    \item $\cos l \cdot \cos m \cdot \cos n = 1$
\end{enumerate}

\item The minimum number of vectors of unequal magnitudes that can add to give a zero resultant is \hfill \textbf{(C)}
\begin{enumerate}[(A)]
    \item 2
    \item 2 (if antiparallel)
    \item 3
    \item 4
\end{enumerate}

\item The angle between $\vec{A} = \hat{i}+\hat{j}-\hat{k}$ and $\vec{B} = 2\hat{i}-\hat{j}+2\hat{k}$ is approximately \hfill \textbf{(B)}
\begin{enumerate}[(A)]
    \item $78.9^\circ$
    \item $101.1^\circ$
    \item $120^\circ$
    \item $150^\circ$
\end{enumerate}

\item In the parallelogram law of addition, if $a$ and $b$ are the magnitudes of two vectors and $\theta$ is the angle between them (tail-to-tail), then the magnitude of their resultant is \hfill \textbf{(D)}
\begin{enumerate}[(A)]
    \item $\sqrt{a^2 + b^2 - 2ab\cos\theta}$
    \item $a + b$
    \item $\sqrt{a^2 - b^2 + 2ab\cos\theta}$
    \item $\sqrt{a^2 + b^2 + 2ab\cos\theta}$
\end{enumerate}

\end{enumerate}

%% ============================================================
\newpage
\topic{Section 4: Graphs}
%% ============================================================

\begin{enumerate}[1.]

\item The slope of the line $y = 4x - 7$ is \hfill \textbf{(C)}
\begin{enumerate}[(A)]
    \item $-7$
    \item $0$
    \item $4$
    \item $-4$
\end{enumerate}

\item For the graph $y = ax^2$, which of the following is correct? \hfill \textbf{(B)}
\begin{enumerate}[(A)]
    \item It is a straight line through the origin
    \item It is a parabola symmetric about the $y$-axis
    \item It is a hyperbola
    \item It passes through $(1, 0)$
\end{enumerate}

\item The graph of $xy = k$ (constant $k$) represents a \hfill \textbf{(D)}
\begin{enumerate}[(A)]
    \item parabola
    \item ellipse
    \item straight line
    \item rectangular hyperbola
\end{enumerate}

\item In an exponential decay graph $y = Ae^{-kx}$ ($k > 0$, $A > 0$), as $x \to \infty$, \hfill \textbf{(A)}
\begin{enumerate}[(A)]
    \item $y \to 0$
    \item $y \to A$
    \item $y \to \infty$
    \item $y \to -A$
\end{enumerate}

\item A straight-line $v$-$t$ graph from $v = 2$ m/s at $t=0$ to $v=10$ m/s at $t=4$ s gives an acceleration of \hfill \textbf{(B)}
\begin{enumerate}[(A)]
    \item $1\,\text{m/s}^2$
    \item $2\,\text{m/s}^2$
    \item $3\,\text{m/s}^2$
    \item $4\,\text{m/s}^2$
\end{enumerate}

\item The displacement in the same $v$-$t$ graph (question 5) between $t=0$ and $t=4$ s equals \hfill \textbf{(C)}
\begin{enumerate}[(A)]
    \item $16$ m
    \item $20$ m
    \item $24$ m
    \item $40$ m
\end{enumerate}

\item The curve $y = 3x^2 - 6x + 5$ has its vertex at \hfill \textbf{(A)}
\begin{enumerate}[(A)]
    \item $(1,\,2)$
    \item $(2,\,1)$
    \item $(0,\,5)$
    \item $(-1,\,14)$
\end{enumerate}

\item A graph of $P$ vs $1/V$ for a fixed mass of ideal gas (Boyle's Law) is \hfill \textbf{(D)}
\begin{enumerate}[(A)]
    \item a parabola through origin
    \item a rectangular hyperbola
    \item an exponential curve
    \item a straight line through the origin
\end{enumerate}

\item The slope of the tangent to the curve $y = f(x)$ at a given point equals \hfill \textbf{(B)}
\begin{enumerate}[(A)]
    \item $\displaystyle\int f(x)\,dx$
    \item $\dfrac{dy}{dx}$ evaluated at that point
    \item the area under the curve at that point
    \item $f(x)$ at that point
\end{enumerate}

\item A graph of $y$ vs $x$ passes through $(1, 3)$ and $(4, 9)$. The area under this line between $x=1$ and $x=4$ is \hfill \textbf{(C)}
\begin{enumerate}[(A)]
    \item $12$ sq. units
    \item $15$ sq. units
    \item $18$ sq. units
    \item $21$ sq. units
\end{enumerate}

\item For exponential graphs, taking $\log$ of both sides converts $y = Ae^{kx}$ to \hfill \textbf{(A)}
\begin{enumerate}[(A)]
    \item a straight line in $\ln y$ vs $x$
    \item a parabola in $y$ vs $\ln x$
    \item a hyperbola in $y$ vs $1/x$
    \item no simpler form
\end{enumerate}

\item The displacement $s(t) = t^3 - 6t^2 + 9t + 2$ has a local maximum at \hfill \textbf{(B)}
\begin{enumerate}[(A)]
    \item $t = 0$
    \item $t = 1$ s
    \item $t = 2$ s
    \item $t = 3$ s
\end{enumerate}

\item The area enclosed by the curve $y = f(x)$ and the $x$-axis between $x=a$ and $x=b$ is given by \hfill \textbf{(D)}
\begin{enumerate}[(A)]
    \item $\dfrac{d}{dx}\displaystyle\int_a^b f(x)\,dx$
    \item the slope of $f$ at $x=a$
    \item $f(b) - f(a)$
    \item $\displaystyle\int_a^b f(x)\,dx$
\end{enumerate}

\item A $v$-$t$ graph gives the displacement by computing \hfill \textbf{(A)}
\begin{enumerate}[(A)]
    \item the area under the $v$-$t$ curve
    \item the slope of the $v$-$t$ curve
    \item the $y$-intercept of the $v$-$t$ curve
    \item the maximum value of $v$
\end{enumerate}

\item For the curve $y = k/x$ with $k > 0$, the $x$- and $y$-axes act as \hfill \textbf{(C)}
\begin{enumerate}[(A)]
    \item tangents to the curve
    \item secants to the curve
    \item asymptotes to the curve
    \item axes of symmetry
\end{enumerate}

\item The graph of $y = 4(1 - e^{-2x})$ for $x \geq 0$ starts at $y = $ \underline{\phantom{XX}} and asymptotically approaches $y = $ \underline{\phantom{XX}}. \hfill \textbf{(B)}
\begin{enumerate}[(A)]
    \item $y=4$ and $y=0$
    \item $y=0$ and $y=4$
    \item $y=-4$ and $y=4$
    \item $y=1$ and $y=4$
\end{enumerate}

\item In a position-time ($s$-$t$) graph, the velocity at any instant is \hfill \textbf{(D)}
\begin{enumerate}[(A)]
    \item the area under the $s$-$t$ curve
    \item the value of $s$ at that instant
    \item the second derivative of $s$
    \item the slope of the tangent to the $s$-$t$ curve at that instant
\end{enumerate}

\item Which of the following graphs would represent the gravitational force $F$ vs distance $r$ between two bodies? \hfill \textbf{(B)}
\begin{enumerate}[(A)]
    \item $F = kr$ (straight line through origin)
    \item $F = k/r^2$ (rectangular hyperbola-like decay)
    \item $F = ke^{-r}$ (exponential decay)
    \item $F = kr^2$ (parabola)
\end{enumerate}

\item Slope $m$ of a straight line in a $y$-$x$ graph satisfies $m = \tan\theta$ where $\theta$ is the angle the line makes with \hfill \textbf{(A)}
\begin{enumerate}[(A)]
    \item the positive $x$-axis
    \item the positive $y$-axis
    \item the origin
    \item the negative $x$-axis
\end{enumerate}

\item A particle's displacement is $s = t^3 - 6t^2 + 9t + 2$. At $t = 2$ s, the acceleration is \hfill \textbf{(C)}
\begin{enumerate}[(A)]
    \item $-6$ m/s$^2$
    \item $2$ m/s$^2$
    \item $0$ m/s$^2$
    \item $6$ m/s$^2$
\end{enumerate}

\item Which pair correctly describes an exponential growth curve? \hfill \textbf{(D)}
\begin{enumerate}[(A)]
    \item $y = A e^{-kx},\; k>0$: steeply rising
    \item $y = Ae^{kx},\; k>0$: rapidly decays to zero
    \item $y = A/x$: grows without bound as $x$ increases
    \item $y = Ae^{kx},\; k>0$: steeply rising without bound
\end{enumerate}

\end{enumerate}

%% ============================================================
\newpage
\topic{Section 5: Differentiation}
%% ============================================================

\begin{enumerate}[1.]

\item $\dfrac{d}{dx}(x^n)$ equals \hfill \textbf{(B)}
\begin{enumerate}[(A)]
    \item $x^{n-1}$
    \item $nx^{n-1}$
    \item $nx^n$
    \item $\dfrac{x^{n+1}}{n+1}$
\end{enumerate}

\item $\dfrac{d}{dx}(\sin x)$ equals \hfill \textbf{(A)}
\begin{enumerate}[(A)]
    \item $\cos x$
    \item $-\cos x$
    \item $\sin x$
    \item $-\sin x$
\end{enumerate}

\item Using the chain rule, $\dfrac{d}{dx}\!\left[e^{-x^2}\right]$ equals \hfill \textbf{(C)}
\begin{enumerate}[(A)]
    \item $e^{-x^2}$
    \item $-e^{-x^2}$
    \item $-2x\,e^{-x^2}$
    \item $2x\,e^{-x^2}$
\end{enumerate}

\item The derivative $\dfrac{d}{dx}(\ln\cos x)$ equals \hfill \textbf{(D)}
\begin{enumerate}[(A)]
    \item $\ln\sin x$
    \item $\cot x$
    \item $\sec x$
    \item $-\tan x$
\end{enumerate}

\item By the product rule, $\dfrac{d}{dx}(x^2\sin x)$ equals \hfill \textbf{(B)}
\begin{enumerate}[(A)]
    \item $2x\cos x$
    \item $2x\sin x + x^2\cos x$
    \item $x^2\cos x$
    \item $2x\sin x - x^2\cos x$
\end{enumerate}

\item By the quotient rule, $\dfrac{d}{dx}\!\left(\dfrac{\sin x}{x}\right)$ equals \hfill \textbf{(A)}
\begin{enumerate}[(A)]
    \item $\dfrac{x\cos x - \sin x}{x^2}$
    \item $\dfrac{x\cos x + \sin x}{x^2}$
    \item $\dfrac{\cos x}{x}$
    \item $\dfrac{-x\cos x + \sin x}{x^2}$
\end{enumerate}

\item The function $f(x) = 2x^3 - 9x^2 + 12x - 5$ has a local maximum at \hfill \textbf{(C)}
\begin{enumerate}[(A)]
    \item $x = 0$
    \item $x = 2$
    \item $x = 1$
    \item $x = 3$
\end{enumerate}

\item If $f''(x) < 0$ at a critical point $x_0$, then $f$ has a \hfill \textbf{(A)}
\begin{enumerate}[(A)]
    \item local maximum at $x_0$
    \item local minimum at $x_0$
    \item saddle point at $x_0$
    \item point of inflection at $x_0$
\end{enumerate}

\item The derivative of $y = x^x$ (for $x>0$) is \hfill \textbf{(D)}
\begin{enumerate}[(A)]
    \item $x^x$
    \item $x \cdot x^{x-1}$
    \item $x^x \ln x$
    \item $x^x(1 + \ln x)$
\end{enumerate}

\item By implicit differentiation of $x^2 + y^2 = r^2$, we get $\dfrac{dy}{dx} =$  \hfill \textbf{(B)}
\begin{enumerate}[(A)]
    \item $\dfrac{y}{x}$
    \item $-\dfrac{x}{y}$
    \item $\dfrac{x}{y}$
    \item $-\dfrac{y}{x}$
\end{enumerate}

\item The charge $q(t) = 3t^2 + 2t + 1$ (coulombs). Current at $t = 3$ s is \hfill \textbf{(C)}
\begin{enumerate}[(A)]
    \item 27 A
    \item 18 A
    \item 20 A
    \item 29 A
\end{enumerate}

\item If $v = u + at$, then $\dfrac{dK}{dt}$ where $K = \tfrac{1}{2}mv^2$ equals \hfill \textbf{(D)}
\begin{enumerate}[(A)]
    \item $mva$
    \item $ma^2t$
    \item $mv^2$
    \item $Fv$ (power delivered)
\end{enumerate}

\item The rate of change $\dfrac{dV}{dt}$ of a sphere's volume when $r = 5$ cm and $\dfrac{dr}{dt} = 2$ cm/s is \hfill \textbf{(A)}
\begin{enumerate}[(A)]
    \item $200\pi$ cm$^3$/s
    \item $100\pi$ cm$^3$/s
    \item $400\pi$ cm$^3$/s
    \item $50\pi$ cm$^3$/s
\end{enumerate}

\item For $f(x) = 2x^3 - 9x^2 + 12x - 5$, the local minimum value is \hfill \textbf{(B)}
\begin{enumerate}[(A)]
    \item 0
    \item $-1$
    \item $-5$
    \item $1$
\end{enumerate}

\item A box with a square base (side $x$) and open top is made from 108 cm$^2$ of material. The height that maximises volume is \hfill \textbf{(C)}
\begin{enumerate}[(A)]
    \item $h = 6$ cm
    \item $h = 2$ cm
    \item $h = 3$ cm
    \item $h = 4$ cm
\end{enumerate}

\item The position $x(t) = t^3 - 12t + 5$. The particle changes direction (for $t \geq 0$) at \hfill \textbf{(A)}
\begin{enumerate}[(A)]
    \item $t = 2$ s
    \item $t = 3$ s
    \item $t = 4$ s
    \item $t = 0$ s
\end{enumerate}

\item $\dfrac{d}{dx}(a^x)$ equals \hfill \textbf{(D)}
\begin{enumerate}[(A)]
    \item $a^x$
    \item $x \cdot a^{x-1}$
    \item $a^x / \ln a$
    \item $a^x \ln a$
\end{enumerate}

\item The electric field $E = kq/r^2$. The value of $\dfrac{dE}{dr}$ is \hfill \textbf{(B)}
\begin{enumerate}[(A)]
    \item $-kq/r^3$
    \item $-2kq/r^3$
    \item $2kq/r^3$
    \item $kq/r^3$
\end{enumerate}

\item The nearest point on the line $y = 2x + 3$ to the origin is \hfill \textbf{(C)}
\begin{enumerate}[(A)]
    \item $(-1,\,1)$
    \item $(0,\,3)$
    \item $\left(-\dfrac{6}{5},\,\dfrac{3}{5}\right)$
    \item $\left(\dfrac{6}{5},\,-\dfrac{3}{5}\right)$
\end{enumerate}

\item At a critical point, the first derivative test shows the function is increasing before and decreasing after the point. This point is a \hfill \textbf{(A)}
\begin{enumerate}[(A)]
    \item local maximum
    \item local minimum
    \item saddle point
    \item global minimum
\end{enumerate}

\item Geometrically, the derivative $\dfrac{df}{dx}$ at a point represents \hfill \textbf{(D)}
\begin{enumerate}[(A)]
    \item the area under $f(x)$ at that point
    \item the average rate of change of $f$
    \item the $y$-intercept of $f$
    \item the slope of the tangent to $y = f(x)$ at that point
\end{enumerate}

\end{enumerate}

%% ============================================================
\newpage
\topic{Section 6: Integration}
%% ============================================================

\begin{enumerate}[1.]

\item $\displaystyle\int x^n\,dx$ (for $n \neq -1$) equals \hfill \textbf{(C)}
\begin{enumerate}[(A)]
    \item $nx^{n-1} + C$
    \item $x^{n-1}/(n-1) + C$
    \item $x^{n+1}/(n+1) + C$
    \item $x^n\ln x + C$
\end{enumerate}

\item $\displaystyle\int \frac{1}{x}\,dx$ equals \hfill \textbf{(B)}
\begin{enumerate}[(A)]
    \item $x + C$
    \item $\ln|x| + C$
    \item $-1/x^2 + C$
    \item $e^x + C$
\end{enumerate}

\item $\displaystyle\int \sin x\,dx$ equals \hfill \textbf{(D)}
\begin{enumerate}[(A)]
    \item $\sin x + C$
    \item $\cos x + C$
    \item $-\sin x + C$
    \item $-\cos x + C$
\end{enumerate}

\item $\displaystyle\int_0^{\pi} \sin x\,dx$ equals \hfill \textbf{(A)}
\begin{enumerate}[(A)]
    \item 2
    \item 0
    \item 1
    \item $\pi$
\end{enumerate}

\item Using substitution $u = x^2$, $\displaystyle\int x\,e^{x^2}\,dx$ equals \hfill \textbf{(C)}
\begin{enumerate}[(A)]
    \item $e^{x^2} + C$
    \item $2e^{x^2} + C$
    \item $\dfrac{e^{x^2}}{2} + C$
    \item $x^2 e^{x^2} + C$
\end{enumerate}

\item By integration by parts, $\displaystyle\int x e^x\,dx$ equals \hfill \textbf{(B)}
\begin{enumerate}[(A)]
    \item $x e^x + C$
    \item $e^x(x-1) + C$
    \item $e^x(x+1) + C$
    \item $x^2 e^x/2 + C$
\end{enumerate}

\item The formula for integration by parts is \hfill \textbf{(D)}
\begin{enumerate}[(A)]
    \item $\displaystyle\int u\,v\,dx = u\displaystyle\int v\,dx$
    \item $\displaystyle\int u\,dv = u\displaystyle\int v\,dx - \displaystyle\int v\,du$
    \item $\displaystyle\int u\,dv = \displaystyle\int v\,du$
    \item $\displaystyle\int u\,dv = uv - \displaystyle\int v\,du$
\end{enumerate}

\item $\displaystyle\int \sin^2 x\,dx$ equals \hfill \textbf{(A)}
\begin{enumerate}[(A)]
    \item $\dfrac{x}{2} - \dfrac{\sin 2x}{4} + C$
    \item $\dfrac{x}{2} + \dfrac{\sin 2x}{4} + C$
    \item $-\dfrac{\cos 2x}{2} + C$
    \item $\dfrac{\cos 2x}{4} + C$
\end{enumerate}

\item $\displaystyle\int_4^{10} \frac{dx}{x}$ equals \hfill \textbf{(C)}
\begin{enumerate}[(A)]
    \item $\ln 4$
    \item $\ln 6$
    \item $\ln\!\left(\dfrac{10}{4}\right) = \ln\!\left(\dfrac{5}{2}\right)$
    \item $\ln 10 - \ln 4 - 1$
\end{enumerate}

\item The physical interpretation of the definite integral $\displaystyle\int_a^b v(t)\,dt$ is \hfill \textbf{(B)}
\begin{enumerate}[(A)]
    \item the acceleration between $t=a$ and $t=b$
    \item the displacement between $t=a$ and $t=b$
    \item the speed at $t=b$
    \item the total distance (regardless of direction)
\end{enumerate}

\item A particle starts from rest with $a(t) = 6t - 2$ m/s$^2$. Its velocity at $t = 3$ s is \hfill \textbf{(D)}
\begin{enumerate}[(A)]
    \item 18 m/s
    \item 15 m/s
    \item 24 m/s
    \item 21 m/s
\end{enumerate}

\item $\displaystyle\int_1^2 (2t-4)\,dt$ equals \hfill \textbf{(A)}
\begin{enumerate}[(A)]
    \item $-3$
    \item $0$
    \item $3$
    \item $-1$
\end{enumerate}

\item $\displaystyle\int e^x \sin x\,dx$ requires integration by parts twice. The result is \hfill \textbf{(C)}
\begin{enumerate}[(A)]
    \item $e^x(\sin x + \cos x)/2 + C$
    \item $e^x\sin x + C$
    \item $e^x(\sin x - \cos x)/2 + C$
    \item $e^x(\cos x - \sin x)/2 + C$
\end{enumerate}

\item $\displaystyle\int \frac{dx}{\sqrt{a^2 - x^2}}$ (using $x = a\sin\theta$) equals \hfill \textbf{(B)}
\begin{enumerate}[(A)]
    \item $\dfrac{1}{a}\tan^{-1}\!\left(\dfrac{x}{a}\right) + C$
    \item $\sin^{-1}\!\left(\dfrac{x}{a}\right) + C$
    \item $\cos^{-1}\!\left(\dfrac{x}{a}\right) + C$
    \item $\dfrac{x}{\sqrt{a^2-x^2}} + C$
\end{enumerate}

\item $\displaystyle\int \frac{dx}{x^2 + a^2}$ equals \hfill \textbf{(D)}
\begin{enumerate}[(A)]
    \item $\dfrac{1}{2a}\ln\!\left(\dfrac{x-a}{x+a}\right) + C$
    \item $\sin^{-1}(x/a) + C$
    \item $a\tan^{-1}(x/a) + C$
    \item $\dfrac{1}{a}\tan^{-1}\!\left(\dfrac{x}{a}\right) + C$
\end{enumerate}

\item The velocity of a particle is $v(t) = 10 - 2t$ m/s. The total distance covered in the first 8 s is \hfill \textbf{(A)}
\begin{enumerate}[(A)]
    \item $34$ m
    \item $16$ m
    \item $25$ m
    \item $50$ m
\end{enumerate}

\item The indefinite integral always includes a constant $+C$ because \hfill \textbf{(C)}
\begin{enumerate}[(A)]
    \item integration is an approximate process
    \item the constant accounts for initial conditions only
    \item the derivative of any constant is zero, so infinitely many antiderivatives exist
    \item integration and differentiation are not inverses of each other
\end{enumerate}

\item $\displaystyle\int_1^e \ln x\,dx$ equals \hfill \textbf{(B)}
\begin{enumerate}[(A)]
    \item $e$
    \item $1$
    \item $e-1$
    \item $0$
\end{enumerate}

\item By the Fundamental Theorem of Calculus, $\dfrac{d}{dx}\displaystyle\int_a^x f(t)\,dt$ equals \hfill \textbf{(D)}
\begin{enumerate}[(A)]
    \item $f(a)$
    \item $F(x) - F(a)$
    \item $\displaystyle\int_a^x f'(t)\,dt$
    \item $f(x)$
\end{enumerate}

\item $\displaystyle\int_0^{\pi/3} \sin x\,dx - \displaystyle\int_{\pi/6}^{\pi/3} \sin x\,dx$ equals $\displaystyle\int_0^{\pi/6} \sin x\,dx$. This follows from the property: \hfill \textbf{(A)}
\begin{enumerate}[(A)]
    \item $\displaystyle\int_a^b f = \displaystyle\int_a^c f + \displaystyle\int_c^b f$
    \item $\displaystyle\int_a^b f = -\displaystyle\int_b^a f$
    \item $\displaystyle\int_a^b kf = k\displaystyle\int_a^b f$
    \item $\displaystyle\int_a^b (f+g) = \displaystyle\int_a^b f + \displaystyle\int_a^b g$
\end{enumerate}

\item The ILATE rule in integration by parts recommends choosing $u$ in the order: \hfill \textbf{(C)}
\begin{enumerate}[(A)]
    \item Exponential, Logarithm, Algebraic, Trig, Inverse Trig
    \item Algebraic, Trig, Logarithm, Inverse Trig, Exponential
    \item Inverse Trig, Logarithm, Algebraic, Trig, Exponential
    \item Trig, Algebraic, Inverse Trig, Logarithm, Exponential
\end{enumerate}

\end{enumerate}

%\vspace{2em}
%\begin{center}
%    \rule{0.5\textwidth}{0.4pt}\\[0.5em]
%    \textit{End of Question Bank} \\
%    \small Total: 22 + 22 + 22 + 20 + 21 + 21 = \textbf{128 Questions}
%\end{center}

\end{document}
